\subsection*{07.05.2026}

\begin{remark}
    Sei $A \in L_b(H)$ positiv für einen Hilbertraum $H$. Es ist $A$ kompakt genau dann wenn $\sigma(A) \setminus \{ 0 \}$ höchstens $0$ als Häufungspunkt hat und
    $$ \dim E(\{ \lambda \}) H < \infty, \quad \forall \lambda \in \sigma(A) \setminus \{ 0 \}. $$
    In dem Fall können wir schreiben
    \[
        A = \int_{\sigma(A) \setminus \{0\}} t \dd E(t) = \sum_{\lambda \in \sigma(A) \setminus \{0\}} \lambda E(\{\lambda\}) = \sum_{n=1}^N \lambda_n \langle \cdot, u_n \rangle u_n,
    \]
    wobei die Reihe sogar gleichmäßig in der Operatornorm konvergiert und $N \in \N \cup \{\infty\}$. Dabei ist $(u_n)$ ein Orthonormalsystem, $\{ \lambda_n \}_{n = 1, \ldots, N} = \sigma(A) \setminus \{0\}$. Weiters können wir die $\lambda_k$ monoton fallend anordnen.
\end{remark}

\begin{definition}
    Seien $H_1, H_2$ Hilberträume. Dann definieren wir
    \[
        S_2(H_1, H_2) = \{ T \in L_b(H_1, H_2) : \sum_{e \in E} \Vert T e \Vert^2 < \infty \},
    \]
    wobei $E \subseteq H_1$ eine (zunächst feste) Orthonormalbasis ist.
\end{definition}

Ist $F \subseteq H_2$ eine Orthonormalbasis, so gilt
\[
    \sum_{e \in E} \Vert T e \Vert^2
    =
    \sum_{e \in E} \sum_{f \in F} \vert \langle Te, f \rangle \vert^2
    =
    \sum_{f \in F} \sum_{e \in E} \vert \langle T^* f, e \rangle \vert^2
    =
    \sum_{f \in F} \Vert T^* f \Vert^2.
\]
Insbesondere ist die rechte Seite unabhängig von der Wahl von $E$, somit ist die obige Definition insbesondere unabhängig von $E$. Außerdem haben wir insbesondere gezeigt $S_2(H_1, H_2)^* = S_2(H_2, H_1)$. Für ein $T \in L_b(H_1, H_2)$ definieren wir nun
\[
    \Vert T \Vert_2^2 = \sum_{e \in E} \Vert Te \Vert^2 \in [0, +\infty].
\]
Wir haben bereits gezeigt $\Vert T \Vert_2 = \Vert T^* \Vert_2$.

Für ein kompaktes $T \in L_b(H_1, H_2)$ ist auch $T^* T$ kompakt als Operator auf $H_1$, deswegen können wir schreiben
\[
    T^*T = \sum_{k=1}^N \lambda_k \langle \cdot, u_k \rangle u_k.
\]
Wir definieren die $s$-Zahlen von $T$ als $s_k(T) \coloneqq \sqrt{\lambda_k}$.

\begin{lemma}
    Sei $T \in L_b(H_1, H_2)$, so ist $T \in S_2(H_1, H_2)$ genau dann, wenn $T$ kompakt ist und
    \[
        \sum_{k=1}^N s_k^2(T) < \infty.
    \]
\end{lemma}

\begin{proof}
    Sei $T \in S_2(H_1, H_2)$ und $\varepsilon > 0$ fest. Dann gibt es eine endliche Teilmenge $M \subseteq E$ mit
    \[
        \sum_{e \in E \setminus M} \Vert Te \Vert^2 < \varepsilon.
    \]
    Für ein $x \in H_1$ mit $\Vert x \Vert < 1$, und $P_M$ der orthogonalen Projektion auf $\spn M$, gilt
    \begin{align*}        
        \Vert (T - T P_M)x \Vert^2
        &=
        \left\Vert \sum_{e \in E} \langle x, e \rangle (T - TP_M) e \right\Vert^2 
        \leq
        \left( \sum_{e \in E} \vert \langle x, e \rangle \vert \, \Vert (T - TP) e \Vert \right)^2 \\
        &\leq
        \left( \sum_{e \in E} \vert \langle x, e \rangle^2 \right) \left( \sum_{e \in E \setminus M} \Vert T e \Vert^2 \right) \leq \Vert x \Vert^2 \varepsilon.
    \end{align*}
    Also ist $\Vert T - TP_M \Vert \leq \varepsilon$, womit wir folgern, dass $T \in K(H_1, H_2)$.

    Ist umgekehrt $T \in K(H_1, H_2)$, so gilt
    \[
        \sum_{k=1}^N s_k^2
        =
        \sum_{k=1}^N \lambda_k
        =
        \sum_{k=1}^N \langle T^* T u_k, u_k \rangle
        =
        \sum_{k=1}^N \Vert T u_k \Vert^2
        =
        \sum_{e \in E} \Vert T e \Vert^2 = \Vert T \Vert_2. \qedhere
    \]
\end{proof}

Insbesondere können wir die $\Vert \cdot \Vert_2$ Norm für kompakte Operatoren also auch über die $s$-Zahlen berechnen.